\makeabstract
    {
    A Comparative Study Investigating Human T-Cell Isolation Kits in Multiple Myeloma
    }
    {
    Morolake Orimoloye\textsuperscript{1}, Alani Perkin\textsuperscript{1,2}, Megan Weivoda\textsuperscript{1,3}
    }
    {
    \textsuperscript{1} Mayo Clinic Graduate School of Biomedical Science, Rochester, MN\\
\textsuperscript{2} Department of Biochemistry and Molecular Biology, Mayo Clinic Graduate School of Biomedical Science, Rochester, MN\\
\textsuperscript{3} Division of Hematology, Mayo Clinic Graduate School of Biomedical Science, Rochester, MN
    }
    {
    T cells are critical components of the immune response against multiple myeloma (MM), with CD3 serving as a marker of T cells and CD8 identifying cytotoxic T-cell subsets capable of targeting tumor cells. Efficient isolation of T cells is therefore important for studying their anti-tumor activity. This study compared the efficiency of the Miltenyi Biotec and PluriSpin negative T-cell isolation kits and evaluated whether differences in isolation methods affected T-cell functionality against MM tumor cells.
    }
    {
    Human T cells were isolated using either the Miltenyi Biotec or PluriSpin negative T-cell isolation kit. Isolation efficiency and CD3+ T-cell purity were assessed using flow cytometry. Isolated patient T cells were subsequently co-cultured with H929 human MM tumor cells to assess cytotoxic T-cell functionality through a tumor-cell killing assay. The killing potential of T cells isolated using each kit was compared.
    }
    {
    Flow cytometric analysis demonstrated greater CD3+ T-cell purity following isolation with the Miltenyi Biotec kit compared with the PluriSpin kit. The Miltenyi isolated T cells also demonstrated greater tumor cell killing potential in the functional killing assay. However, T cells isolated using the PluriSpin kit retained the ability to kill MM tumor cells, indicating that both isolation methods yielded functionally active T cells. Visual assessment of the killing assays did not reveal large differences between the two isolation methods.
    }
    {
    The Miltenyi Biotec negative T-cell isolation kit provided improved CD3+ T-cell isolation compared with the PluriSpin kit based on flow cytometric purity assessment. T cells isolated using the Miltenyi kit demonstrated greater cytotoxic potential against MM tumor cells, although PluriSpin-isolated T cells also maintained tumor-cell killing activity. These findings suggest that the choice of T-cell isolation method may influence both the purity and functional assessment of T cells in studies of anti-tumor immunity in multiple myeloma.
    }

    \newpage
    